\section{实验三：机械臂抓取与搬运}

\subsection{实验目的}

本实验围绕机械臂从一点抓取物体并搬运到另一点的任务编写动作规划与仿真程序。与单纯的末端轨迹控制相比，抓取搬运程序需要同时处理关键位姿序列、夹爪开合、桌面安全高度和未来实机执行保护。

\subsection{任务流程设计}

程序位于 \path{daran/task2/pick_and_place.py}，采用参考位姿回放方式规划 A 点抓取和 B 点放置。已有文件 \path{ref_pose.json} 中记录了四个关键点：

\begin{table}[H]
  \centering
  \caption{抓取搬运实验示教点}
  \label{tab:pick-place-points}
  \small
  \begin{tabularx}{0.82\textwidth}{L{1.5cm}L{2.5cm}Y}
    \toprule
    顺序 & 名称 & 含义 \\
    \midrule
    1 & A\_safe & A 点上方安全位 \\
    2 & A\_grasp & A 点抓取位 \\
    3 & B\_safe & B 点上方安全位 \\
    4 & B\_place & B 点放置位 \\
    \bottomrule
  \end{tabularx}
\end{table}

完整执行序列为：
\[
\text{park}\rightarrow A_{\text{safe}}\rightarrow A_{\text{grasp}}\rightarrow
A_{\text{safe}}\rightarrow B_{\text{safe}}\rightarrow B_{\text{place}}\rightarrow
B_{\text{safe}}\rightarrow \text{park}.
\]
其中到达 A 点抓取位后闭合夹爪，到达 B 点放置位后打开夹爪。

该流程采用“上方安全点”和“接触操作点”分离的设计。真实机械臂在桌面附近水平移动时，夹爪或物体可能与桌面、障碍物发生干涉，因此程序规划末端先移动到目标点上方，再沿近似竖直方向下降到抓取或放置位。这个设计为后续实机实验预留了安全动作模板，但其实际避碰效果尚未验证。

\subsection{示教点记录}

\path{ref_pose.json} 中的四组参考数据包含六个机械臂关节角和夹爪角度。本组将其作为仿真输入，没有通过真实机械臂重新示教或验证。表 \ref{tab:pick-place-joints} 给出了四个关键点的模型关节角，表 \ref{tab:pick-place-fk} 给出了由 DH 正运动学计算得到的末端坐标。

\begin{table}[H]
  \centering
  \caption{抓取搬运示教点关节角}
  \label{tab:pick-place-joints}
  \small
  \begin{tabularx}{\textwidth}{L{1.6cm}YYYYYY}
    \toprule
    点位 & $q_1$ & $q_2$ & $q_3$ & $q_4$ & $q_5$ & $q_6$ \\
    \midrule
    A\_safe & 22.5 & 58.9 & -46.9 & -1.9 & 84.5 & 22.7 \\
    A\_grasp & 17.9 & 53.4 & -53.7 & 7.3 & 92.2 & 12.0 \\
    B\_safe & -24.4 & 60.9 & -50.7 & -0.6 & 83.6 & -24.8 \\
    B\_place & -23.8 & 55.6 & -57.7 & 7.2 & 82.7 & -24.9 \\
    \bottomrule
  \end{tabularx}
\end{table}

\begin{table}[H]
  \centering
  \caption{示教点正运动学末端坐标}
  \label{tab:pick-place-fk}
  \small
  \begin{tabularx}{0.78\textwidth}{L{1.8cm}YYY}
    \toprule
    点位 & $x$/mm & $y$/mm & $z$/mm \\
    \midrule
    A\_safe & 254.5 & 161.5 & 139.2 \\
    A\_grasp & 278.6 & 149.1 & 95.2 \\
    B\_safe & 288.1 & -74.4 & 136.6 \\
    B\_place & 299.8 & -76.8 & 91.7 \\
    \bottomrule
  \end{tabularx}
\end{table}

从表 \ref{tab:pick-place-fk} 可以看到，A 点安全位到抓取位的末端高度下降约 \SI{44.0}{mm}，B 点安全位到放置位的高度下降约 \SI{44.9}{mm}，均小于程序设定的 \SI{120}{mm} 下降幅度上限。同时，两个 safe 点高度均高于 \SI{130}{mm}，满足脚本中的安全高度要求。这说明示教点并非只记录目标位置，还记录了完整任务中“安全接近”和“低位接触”的空间关系。

\subsection{关键位姿与夹爪控制}

脚本定义了固定 park 位姿 $[0,90,0,0,0,0]^\circ$ 作为任务起止安全姿态。示教点中的关节角全部经过限位检查，A/B 安全点还需满足末端高度不低于 \SI{130}{mm}，safe 点到 grasp/place 点的下降幅度不得超过 \SI{120}{mm}。

夹爪控制代码未直接使用原始 \texttt{grasp()} 接口，而是根据程序中给定的换算关系调用底层 \path{set_angle_adaptive}：
\[
\theta_g = w \cdot \frac{-180}{\pi d},
\]
其中 $w$ 为夹爪开口宽度，$d=\SI{10}{mm}$ 为夹爪传动直径。程序设置的张开宽度为 \SI{6.0}{mm}，夹紧宽度为 \SI{1.8}{mm}，闭爪力参数约为 \SI{40}{N}；这些参数未在实物夹取中验证。

\subsection{实验步骤与安全检查}

抓取搬运实验首先通过示教或已有文件获得 A、B 两处关键位姿，然后由程序读取 \path{ref_pose.json} 并进行格式检查。随后程序根据 DH 模型计算各示教点的末端高度，确认 safe 点高于桌面安全线，且从 safe 点下降到 grasp/place 点的距离处于允许范围内。只有在这些检查通过后，程序才会生成完整动作序列。

仿真阶段主要观察三类现象：第一，机械臂是否能从 park 姿态平滑到达各关键点；第二，下降和抬升动作是否基本沿竖直方向进行；第三，关节角是否存在突变或接近限位的情况。实机执行阶段则需要额外确认夹爪开口是否与物体尺寸匹配，目标物体是否位于示教抓取点附近，以及抓取后物体是否会因重心偏移而滑落。

程序将 10 个任务步骤组织为统一计划，其中 7 个为机械臂移动段，3 个为夹爪动作段。每个移动段使用五次多项式关节插值生成 40 个采样点；每个夹爪动作在仿真中用短暂停顿表示，以便观察动作时序。因此仿真轨迹不仅展示机械臂从 A 点移动到 B 点的空间路径，也展示了抓取前张爪、抓取后闭爪和放置后张爪的时序位置。

安全检查分为示教点检查和整条轨迹检查两类。示教点检查关注四个关键点本身是否合理，例如关节角是否越界、safe 点是否高于抓取点、下降距离是否过大；整条轨迹检查则关注插值过程中是否出现中间姿态低于桌面或末端高度异常。这样可以避免一种常见问题：起点和终点都安全，但两点之间的插值姿态穿过危险区域。

\subsection{程序实现}

抓取搬运程序沿用了实验二中的 DH 模型和关节限位。主要函数包括：

\begin{itemize}
  \item \texttt{load\_waypoints()}：读取示教点并检查数据长度和关节限位。
  \item \texttt{safety\_check\_reference\_set()}：计算示教点正运动学位置，检查桌面高度、安全点高度和下降幅度。
  \item \texttt{build\_plan()}：构造移动和夹爪动作序列。
  \item \texttt{verify\_ee\_z\_floor()}：检查整条轨迹末端高度下限。
  \item \texttt{MOVE\_REAL\_ROBOT}：控制是否真实下发串口指令。
\end{itemize}

值得注意的是，脚本头部注释要求先将 \texttt{MOVE\_REAL\_ROBOT} 设为 \texttt{False} 生成仿真动画，再由操作者决定是否进入实机模式。本组始终保持仿真模式，没有将动作下发到真实机械臂。

为未来实机执行，代码没有直接使用库中无超时的 \texttt{pose\_done()} 或 \texttt{grasp\_done()} 轮询，而是实现了带超时的等待函数。其设计目标是在关节未按时到位时中止任务，并区分张爪故障与闭爪接触物体两类情况。由于本组未连接硬件，这些超时与接触判断逻辑仅完成代码检查，尚未验证实际行为。

抓取搬运的核心动作序列和夹爪宽度到角度的换算如代码 \ref{lst:exp3-code} 所示。代码中将移动动作和夹爪动作放在同一计划列表中，便于仿真和实机执行共用同一套流程。

\lstinputlisting[
  language=Python,
  caption={抓取搬运动作序列与夹爪宽度换算},
  label={lst:exp3-code}
]{code/exp3_pick_and_place.py}

\subsection{仿真结果与分析}

图 \ref{fig:task3-pick-place} 为抓取搬运动画中抽取的关键帧。仿真模型先回到竖直 park 姿态，再移动到 A 点上方；下降到抓取位后显示闭爪时序，随后抬升、平移至 B 点上方，再下降并显示开爪时序。该动画用于检查动作顺序和模型姿态，不代表夹爪已经接触或搬运真实物体。

\begin{figure}[H]
  \centering
  \includegraphics[width=0.82\textwidth]{task3_pick_place_keyframes.jpg}
  \caption{抓取搬运过程仿真关键帧}
  \label{fig:task3-pick-place}
\end{figure}

关键帧中可以看到，A 点和 B 点均设置 safe/grasp 或 safe/place 两级高度，模型在低位只执行短距离接近和离开动作，主要水平移动发生在安全高度附近。这与表 \ref{tab:pick-place-fk} 中 safe 点高于 grasp/place 点约 \SI{44}{mm} 的数据一致，也说明参考点不是孤立位置，而是服务于完整任务时序的空间约束。

本实验的精度主要依赖示教点质量。若 A\_grasp 点偏离物体中心，夹爪可能只夹住物体边缘；若 B\_place 点高度过低，放置时可能挤压桌面；若 safe 点高度不足，平移过程中可能带动物体碰撞其它物体。因此示教点不仅记录位置，还隐含了操作者对安全高度、抓取方向和物体尺寸的经验判断。后续视觉抓取的目标是将其中“抓取点位置”由视觉自动给出，但 safe 点构造和夹爪动作顺序仍沿用本实验的思想。

与实验二相比，本实验的核心难点不再是单条直线是否足够直，而是任务序列是否满足接触操作的约束。抓取动作必须在物体上方张爪，在接触高度闭爪，再抬升到安全高度后水平搬运；放置动作则必须在目标上方下降，松开夹爪后再抬升离开。任何一步时序错误都可能导致空抓、碰撞或拖动物体。因此本实验从运动控制上升到了任务规划层面。

从后续实机风险看，抓取搬运既包含运动学误差，也包含夹爪和物体接触误差。即使关节到达参考角度，如果物体位置与参考抓取点存在偏差，夹爪仍可能抓偏；夹爪力过小可能导致滑落，过大则可能挤压物体。任务 2 中的夹爪宽度、闭爪力和停顿时间目前只是程序参数，需通过实物实验重新调整。

\subsection{失败模式与保护策略}

抓取搬运任务比单纯运动控制更容易失败，因为它涉及机械臂、夹爪、物体和桌面之间的接触。根据程序结构可以预判四类失败模式。第一类是参考点错误，例如 A\_grasp 没有对准物体中心，或 B\_place 过低；第二类是轨迹错误，例如 safe 点高度不足；第三类是夹爪错误，例如张爪不到位或闭爪力不足；第四类是通信和执行错误，例如关节未按时到位或串口读取失败。这些模式尚未通过本组实物实验观察。

脚本中的保护策略正是围绕这些失败模式展开。参考点错误通过关节限位、safe 高度和下降幅度检查发现；轨迹错误通过整条插值轨迹的桌面高度扫描发现；夹爪与执行错误则预期通过宽度反馈轮询、关节到位轮询、超时和异常中止处理。这些保护逻辑是否能够覆盖真实机械臂故障仍需后续测试。

需要注意的是，保护策略只能降低风险，不能消除所有风险。例如程序可以检查末端坐标系高度，但无法直接知道夹爪外形、物体实际高度和桌面上其它障碍物的位置；程序可以等待关节到位，但无法完全判断物体是否已经被稳定夹住。因此实机执行前仍然需要操作者观察仿真动画、清理工作区、确认物体位置，并在任务运行时保持可中止状态。

\subsection{小结}

本实验在代码和仿真中将前一实验的轨迹规划扩展为包含夹爪动作的任务级流程。其关键不是单个点的运动，而是多个安全位姿之间的顺序约束。后续视觉抓取程序也是先通过视觉确定抓取点，再生成类似的 safe-grasp-safe 序列，因此本实验为视觉抓取提供了动作模板。

本实验也暴露了示教回放方法的局限性：它要求物体每次都放在相同位置，且放置点环境不能发生明显变化。一旦物体初始位置偏移，固定示教点就无法自动修正。因此实验四和实验五继续引入相机标定与视觉定位，使系统能够由图像实时确定目标位置，从而把本实验中的固定 A\_grasp 点替换为视觉计算得到的抓取点。
